% Generated draft: requires mathematical and source review.
\documentclass[11pt]{article}
\usepackage[margin=25mm]{geometry}
\usepackage{fontspec}
\setmainfont{Latin Modern Roman}
\usepackage{amsmath,amssymb,mathtools,unicode-math}
\setmathfont{Latin Modern Math}
\usepackage{enumitem,needspace,xurl,hyperref}
\setlength{\parindent}{0pt}
\setlength{\parskip}{5pt}
\setlength{\emergencystretch}{3em}
\newcommand{\N}{\mathbb{N}}
\newcommand{\Z}{\mathbb{Z}}
\newcommand{\Q}{\mathbb{Q}}
\newcommand{\R}{\mathbb{R}}
\newcommand{\C}{\mathbb{C}}
\newcommand{\F}{\mathbb{F}}
\newcommand{\A}{\mathbb{A}}
\newcommand{\PP}{\mathbb{P}}
\newcommand{\E}{\mathbb{E}}
\newcommand{\eps}{\varepsilon}
\newcommand{\dd}{\,\mathrm{d}}
\begin{document}
\section*{UnsolvedMath ID: 30001170}
\textbf{Source problem code:} catalogue entry
\subsection{Definitions and mathematical statement}
The supplied record collects several questions about Schrödinger spectral gaps. It does not specify the operator realization, boundary conditions, admissible potential class, or every target inequality. The notation below therefore gives a conditional framework, rather than supplying these missing hypotheses.

Let \(H_{\Omega,V}\) denote a self-adjoint realization of the differential expression
\[
-\Delta+V
\]
on a domain \(\Omega\), whenever the chosen realization has discrete eigenvalues
\[
\lambda_1(\Omega,V)\leq\lambda_2(\Omega,V)\leq\cdots,
\]
listed with multiplicity. Define its fundamental gap by
\[
\Gamma(\Omega,V)=\lambda_2(\Omega,V)-\lambda_1(\Omega,V).
\]
In the one-dimensional questions, \(\Omega=I\) is an interval; when a finite diameter is used, write \(d=\operatorname{diam}(\Omega)>0\). The source record does not fix boundary conditions or the regularity and integrability assumptions needed to define \(H_{\Omega,V}\).

For a nonzero real eigenfunction \(u_j\) corresponding to \(\lambda_j\), interpret the normalized potential mean, provided the integrals exist, as
\[
M_j(V)=\langle Vu_j^2\rangle
:=\frac{\displaystyle\int_\Omega V(x)u_j(x)^2\,\dd x}
{\displaystyle\int_\Omega u_j(x)^2\,\dd x}.
\]
This is an explicit interpretation of the supplied notation; its normalization is not defined in the record. For complex eigenfunctions, the corresponding expectation would use \(|u_j|^2\), but that convention is not specified by the source record. Choices within multiple eigenspaces would also need to be addressed.

\begin{enumerate}
\item \textbf{One-dimensional potential classes.} Characterize the largest class of potentials for which the intended fundamental-gap inequality holds. More formally, after fixing an ambient admissible class \(\mathcal A\), the operator realization on \(I\), and the intended lower bound \(B(d)\), the target class would be
\[
\mathcal C_{I,B}
=\{V\in\mathcal A:\Gamma(I,V)\geq B(d)\}.
\]
The record does not specify \(\mathcal A\), \(B\), or whether the estimate is required on one interval or uniformly over intervals. Thus this expression identifies the missing data rather than selecting a particular inequality.

\item \textbf{Potential means.} Determine whether the intended hypotheses imply
\[
M_{j+1}(V)\geq M_j(V)\qquad\text{for every }j\geq1.
\]
Here nondecreasing is one possible meaning of ``increasing''; whether strict inequalities are intended, and for which potentials and boundary conditions, remains unresolved.

\item \textbf{Potential-dependent strengthening.} Determine whether the intended gap inequality admits a stronger lower bound containing an explicitly specified correction involving \(V\). Schematically, such a question would ask for
\[
\Gamma(\Omega,V)\geq B(d)+R(\Omega,V),
\]
where the correction yields an improvement in the intended class. Neither the functional form nor the sign, normalization, or applicable hypotheses for \(R\) are supplied.

\item \textbf{Higher-dimensional methods and bounds.} Determine whether Lavine's method extends to higher dimensions, and whether Brascamp--Lieb or Korevaar methods give analogous or stronger spectral-gap estimates. The original statement specifically asks whether these methods improve a bound written as \(\pi^2/d^2\). The record does not state the hypotheses of that bound, the desired improvement, or its precise relationship to the fundamental-gap estimate in the other questions.

\item \textbf{Zero-potential minimizing domains and degeneration.} The original statement also asks, for \(V\equiv0\) on convex domains of fixed diameter, whether the infimum of the fundamental gap is not attained and whether minimizing domains shrink to a segment, with the historical aim of completing the gap conjecture. Conditionally, if \(\mathcal K_{n,d}\) denotes the intended class of admissible convex domains in \(\R^n\) with diameter \(d\), define
\[
g_{n,d}=\inf_{\Omega\in\mathcal K_{n,d}}\Gamma(\Omega,0).
\]
The questions concern nonattainment of \(g_{n,d}\) and degeneration to a segment for domains \(\Omega_k\) satisfying
\[
\Gamma(\Omega_k,0)\longrightarrow g_{n,d}.
\]
The dimension range, boundary conditions, admissible domain regularity, convergence topology, allowed rigid motions or subsequences, and whether the assertion concerns every minimizing sequence are not specified. Consequently, no precise universal degeneration assertion is determined by the supplied text.
\end{enumerate}

The supplied dated curation labels the collection partially solved and reports that the fundamental gap conjecture for convex domains and convex potentials was proved. It also says that listed magnetic, higher-dimensional, and strengthening questions are not all settled by that theorem. No magnetic operator or magnetic question is actually formulated in this record. These are supplied status qualifications, not independently verified mathematical conclusions, and the historical completion question above is not asserted here to remain open.
\subsection{Short English statement}
This collection asks which one-dimensional potentials satisfy a spectral-gap bound, whether normalized potential expectations increase with the eigenvalue index, and whether potential-dependent corrections or higher-dimensional methods improve the estimate. The original record also asks about nonattainment and degeneration to a segment for zero-potential convex domains of fixed diameter. Essential operator hypotheses and several target bounds are missing. The supplied partially solved label and literature assessment are dated curation; this statement-only draft requires human mathematical review.
\subsection{Sources}
\begin{itemize}
\item Local UnsolvedMath registry: supplied complete record for permanent ID \(30001170\), entitled \emph{Refinements and Degenerations in Schrödinger Spectral Gaps}; exact source problem code \texttt{catalogue entry}. Collection: Oberwolfach Reports Open Problems; proposed year: 2009. No individual proposer is supplied.
\item Supplied primary report citation: \emph{Low Eigenvalues of Laplace and Schrödinger Operators} (2009), Oberwolfach Reports, pp.\ 355--428, DOI \texttt{10.4171/owr/2009/06}. Supplied source URL and literature-source link: \url{https://doi.org/10.4171/owr/2009/06}. The record describes the report as documenting workshop discussions and open problems; no exact internal page or speaker attribution for this entry is supplied.
\item Supplied literature reference, labeled ``Andrews--Clutterbuck fundamental gap theorem'': \url{https://doi.org/10.1090/S0894-0347-2011-00699-1}.
\item Literature assessment date: \texttt{2026-08-21}. This is supplied dated curation, not independently verified. Its label is ``partially solved''; it reports the fundamental gap conjecture for convex domains and convex potentials as proved, while saying the listed magnetic, higher-dimensional, and strengthening questions are not all settled by that theorem. The source itself requests independent verification of time-sensitive status and cited claims.
\item The registry labels the statement \texttt{corrected\_verified} and reports that its mathematical content was checked against the official Oberwolfach report, with notation, context, and question form normalized. That verification claim is supplied provenance and has not been independently checked for this draft.
\end{itemize}
\end{document}
